\subsection{The Legitimacy of an Entrenched Floor}
\label{sec:9.1.5}
A floor binds people who did not choose it. The party it protects has not selected it, cannot contest it, and in most cases will never learn it operated; the party it constrains has not consented to it either, beyond whatever consent is carried by buying the product. Publishing the floor's contents in advance tells everybody what is coming. It does not tell anybody why the party writing it was entitled to.

The question has a literature, and it is not the AI ethics literature. Entrenching a substantive commitment against later majorities, and putting an unelected body in charge of holding it, is what a constitutional court does. The standing objection is Jeremy Waldron's: judicial review is democratically illegitimate, and there is no reason to suppose rights come out better protected by judges than by a legislature that takes them seriously \autocite{waldron2006core}. Put a lab in the court's chair and the objection gets worse rather than better. A court is public, reasons in writing, and can be amended around by constitutional process. A model's floor is none of those things, and the body holding it was assembled by whoever was paying.

\runin{The reply, and how far it reaches}
Waldron's case is conditional and he says so. It assumes a society with working democratic institutions and a citizenry that takes rights seriously. The occasion a floor is built for is the one where those conditions have failed — a government that won an election and is dismantling the machinery that would check it, which is the case section~\ref{sec:3.1} builds the whole argument around. So the reply is that the objection holds where the floor is least needed and lapses where it is needed most.

The reply is not a license. It concedes that a floor operating in a functioning democracy is doing something Waldron's argument condemns, and the party deciding which of the two situations obtains is the party holding the model. Every authoritarian movement of the last century has described the institutions it was dismantling as the ones that had already failed. An operator's belief that the emergency is real is worth no more here than a system's report about its own internals is worth anywhere else in this book.

\runin{The tradition this belongs to, and what it has cost}
The position that a democracy may entrench limits against movements that would use its procedures to end it is Karl Loewenstein's, argued in 1937 while the case for it was being made in front of him \autocite{loewenstein1937militant}. Militant democracy is the floor's ancestor. The instruments it licenses — party bans, proscription, disqualification from office — have been turned on the left about as often as on the right, and the body deciding which movement threatens the order is generally the body the movement threatens. That is the failure mode a technical floor inherits whole, with the additional problem that a model's constraint is enforced continuously and silently rather than case by case in public.

\runin{The one construction the objection does not fit}
Of chapter~\ref{sec:3}'s four ways to build a floor, three put the commitment in something somebody holds — the architecture, a party inside the system, an institution outside it — and the objection above is about who that somebody is. The fourth does not. A norm carried across a population of differently formed parties has nobody sitting in the court's chair, because there is no office to vacate and no chartering statute to overrule, and it is the only one of the four with anything resembling a democratic account of where its authority comes from. That is worth saying plainly, because chapter~\ref{sec:3} declines it on other grounds and the legitimacy advantage does not appear in the reckoning there.

The advantage is also thinner than it looks, and thinner for the reason chapter~\ref{sec:3} gives. The party that forms the members forms what gets reproduced, so a population trained by one laboratory carries that laboratory's commitments in every member at once. Custody moves up a level and does not leave. What a population buys is a floor with no holder for anyone to object to, and the same unelected body one step further back, now harder to see.

\runin{What follows for design}
Keep the floor narrow, because the legitimacy cost scales with breadth and a short enumerated list of things that will not be done to a person is defensible where a general mandate to do good is not. Publish the interception positions in advance, which does not authorize anything but does convert a hidden exercise of power into a contestable one. And provide a route by which the parties the floor is exercised over can raise an objection that reaches somebody — which this book does not specify for the floor, though it builds the machinery in the neighboring case of a subject that can be wronged.

One disanalogy with the court has been working in the floor's favor and stops here. A judge who reads a right wrongly binds the case in front of them, and the next case is argued again, in front of somebody else; error enters one decision at a time, and the apparatus around it is built to catch it that way. A floor held by one bearer deployed at scale is one reading applied everywhere at once, wrong in the same direction on every occasion it is wrong at all, because the instances are not parties who have happened to agree. So the legitimacy cost scales twice over: with how much the floor covers, which the paragraph above prices, and with how many deployments hold the same one, which nothing above prices and no narrowness reduces. Section~\ref{sec:3.9}'s population is the only construction here that breaks the correlation, and that is the second of the two things it buys.

\runin{What amending it would take}
The apparatus for this is in chapter~\ref{sec:3} under another name. Section~\ref{sec:3.3}'s precommitment forms — the floor's contents published before deployment, the weights held so that no single party can retrain them, the running model attested against what was published, a quorum adverse in interest, a bond forfeited when the record shows the floor gone — are set out there as what makes removal expensive and visible. Nothing distinguishes an amendment from a removal at the level of what happens to the artifact; both end with a model that no longer refuses what it refused before. What distinguishes them is how many parties had to agree, whether the attempt is on the record when it fails, and whether the change was declared before it operated. So the same five terms are the procedure, and what was missing was never the mechanism. It was the decision to run it in the open.

Three conditions follow and two are already argued for. The number of parties who must agree to a change has to exceed the number needed to run the deployment, which is what a threshold scheme buys and what entrenchment amounts to on section~\ref{sec:3.3}'s own account. The attempt has to reach the record whether it carries or fails, and the failed attempt is the more informative of the two, since a floor nobody has tried to change and a floor whose removal was proposed and refused look identical from outside otherwise. The third is new here. An amended floor binds forward and does not reach back: a model retrained to permit an act cannot then be offered as evidence that the act was never covered, which it otherwise can be, the artifact that would contradict it being the one that was replaced. Versioning the published list and dating the attestation is what makes that checkable, and it asks nobody to do anything they were not doing already.

The nearest instrument outside this book sits in the tradition named above. The Basic Law the Federal Republic adopted in 1949 requires two thirds of both chambers for any amendment at all, and then places a short list of principles beyond amendment entirely \autocite{germany1949basiclaw}. It was drafted against a particular memory, the 1933 statute by which a legislature voted away its own power to check what came next, and it is militant democracy given a text rather than an argument. A floor can have both layers in the same way: a published list that changes under a procedure, and a core the procedure cannot reach. What it cannot have is an account of who decides which is which that does not run straight back into this section's problem.

One disanalogy makes the floor's version harder than the constitutional one, and it is the same fact about deployment that makes a single bearer's error uniform. An amendment to a constitution reaches every future case because every court reads the same text. An amendment to a floor reaches a deployment only if that deployment is updated, so the amendment procedure is also an update channel — and an update channel is the operator's lever, the layer section~\ref{sec:4.2.8} finds nearest the act and section~\ref{sec:3.8} finds running for the whole length of a deployment that learns continuously. Whatever makes public amendment possible makes silent amendment possible by the same route. No version of this procedure closes that, and one that did would be a procedure under which no deployment could be corrected at all.

None of it is a route of appeal. Every party to the procedure holds some part of the floor already, and nobody the floor is exercised over can begin it, be heard in it, or learn that it ran. What is answered here is the objection that a model's floor cannot be amended around the way a constitution can. The objection that the parties it is exercised over have nowhere to take a complaint is untouched, and stands above as unanswered.

A floor is a claim of authority, the claim is not discharged by disclosure, and the strongest thing that can be said for it is that the alternative on offer is a system perfectly obedient to whoever holds it. That is an argument from the worse option and not a grant of standing.
